% GENERADO desde borrador_reescrito/02_marco_teorico.md — no editar a mano

\section{Marco teórico}

\subsection{Marco conceptual}

Esta sección define los cuatro conceptos sobre los que se apoya el estudio: el trámite, la carga administrativa, la interoperabilidad y el principio de \enquote{solo una vez}. Se presentan en ese orden porque cada uno se apoya en el anterior: el trámite es el espacio en el que el ciudadano entrega información a la entidad; la carga administrativa describe el costo que esa entrega le supone; la interoperabilidad es el medio por el cual la entidad puede obtener la información sin pedirla, y el principio de \enquote{solo una vez} es el criterio que orienta ese cambio.

\subsubsection{Trámite}

En Colombia, el Decreto 088 de 2022 define el trámite como el \enquote{conjunto de requisitos, pasos o acciones, regulados por el Estado, dentro de un proceso misional, que deben efectuar los ciudadanos, usuarios o grupos de interés} ante una entidad pública o un particular que cumple funciones públicas, para hacer efectivo un derecho, ejercer una actividad o cumplir una obligación prevista o autorizada por la ley \parencite[art.~2.2.20.3, num.~13]{colombia2022decreto088}. Una definición cercana se encuentra en el estudio regional de \textcite{roseth2018complicada}, para quienes los trámites son el conjunto de requisitos, pasos o acciones por los cuales los individuos o las empresas \enquote{piden o entregan información a una entidad pública} con el fin de obtener un derecho o cumplir una obligación (p.~36); los autores indican que adaptaron esa definición de las de la Comisión Federal de Mejora Regulatoria de México y del \enquote{Departamento de Función Pública de Colombia} \parencite[p.~36, nota~1]{roseth2018complicada}. Ambas definiciones incluyen los requisitos como parte del trámite. A partir de ellas, en este estudio los documentos que la entidad exige al ciudadano se entienden como requisitos mediante los cuales este entrega información para obtener el derecho o cumplir la obligación.

Para analizar esos requisitos, el estudio los clasifica según quién expide o custodia el documento, con base en tres distinciones que hace la normativa colombiana. La primera agrupa los documentos que reposan en la propia entidad ante la cual se adelanta el trámite, que el Decreto Ley 019 de 2012 prohíbe exigir \parencite[art.~9]{colombia2012decreto019}. La segunda agrupa los documentos de competencia de otras autoridades, cuya presentación la Ley 962 de 2005 prohíbe solicitar \parencite[art.~1, num.~1]{colombia2005ley962}; dentro de ella se distinguen los que constan en registros públicos, que, mientras sus responsables no se integren al servicio ciudadano digital de interoperabilidad, deben poder consultarse en línea y de forma gratuita, de modo que el ciudadano queda eximido de aportar el certificado \parencite[art.~10, par.~3]{colombia2019decreto2106}. La tercera es una categoría residual, propuesta por este estudio, que reúne la información que solo el ciudadano puede aportar, como sus propias declaraciones o los formularios que diligencia, y que ninguna base de datos pública contiene. Además, \textcite{roseth2018complicada} señalan que los trámites \enquote{implican costos de transacción (algunos más que otros)} (p.~36) y reportan que, según los encuestados de Latinobarómetro en 2017, realizar su último trámite les tomó en promedio aproximadamente 5,4 horas en América Latina y el Caribe (p.~47). Ese costo es el objeto del concepto que se presenta a continuación.

\subsubsection{Carga administrativa}

Los costos que asume una persona al interactuar con el Estado se estudian bajo el concepto de carga administrativa. \textcite{moynihan2015administrative} la conceptualizan como una función de los costos de aprendizaje, psicológicos y de cumplimiento que experimentan los ciudadanos en sus interacciones con el Estado, y sostienen que el nivel de esa carga y su distribución entre el Estado y el individuo suelen ser resultado de decisiones políticas deliberadas y no simplemente de accidentes históricos o descuidos (p.~43). De los tres tipos de costo, los de cumplimiento son las cargas de seguir las reglas y requisitos administrativos; para quien solicita un servicio, son los costos de diligenciar formularios o de aportar documentos que acrediten su condición \parencite[p.~46]{moynihan2015administrative}. Trabajos posteriores de los mismos autores mantienen esa tipología y ubican el papeleo y la documentación dentro de los costos de cumplimiento \parencite{herd2025administrative} \pendiente{CONFIRMAR}. Este estudio no mide la carga que experimentan los ciudadanos del Huila; usa el concepto para interpretar cada documento que se exige como un costo de cumplimiento y, en consecuencia, dejar de exigirlo como una forma de trasladar esa parte de la carga a la administración. En ese sentido, la reducción de la carga administrativa es el beneficio esperado para el ciudadano que justifica la propuesta. El concepto siguiente describe el medio que hace posible ese traslado.

\subsubsection{Interoperabilidad}

El Marco de interoperabilidad para Gobierno Digital, publicado por el Ministerio de Tecnologías de la Información y las Comunicaciones (MinTIC), acoge como definición de interoperabilidad \enquote{la capacidad de las organizaciones para intercambiar información y conocimiento en el marco de sus procesos de negocio} \parencite[p.~33]{mintic2019marco}; según esa definición, el propósito del intercambio es que ciudadanos, empresas y otras entidades puedan recibir servicios en línea gracias a que los sistemas de las organizaciones comparten datos. El mismo documento describe que, en un escenario funcional de interoperabilidad, cuando una entidad requiera comprobar una circunstancia que obra en otra entidad, \enquote{no debe solicitarle dicha información al usuario, sino que deberá obtenerla de la entidad respectiva} por un mecanismo digital \parencite[p.~19]{mintic2019marco}. A ese mecanismo el marco lo denomina servicio de intercambio de información, y lo define como un recurso tecnológico que, mediante un conjunto de protocolos y estándares, permite el intercambio de información \parencite[p.~10]{mintic2019marco}. En la normativa colombiana, los Servicios Ciudadanos Digitales se definen como el conjunto de soluciones y procesos transversales que brindan al Estado capacidades para su transformación digital y para una adecuada interacción con el ciudadano \parencite[art.~2.2.17.1.4, num.~13]{colombia2020decreto620}, y entre los servicios base se incluye el servicio de interoperabilidad, que brinda las capacidades necesarias para garantizar el adecuado flujo de información entre los sistemas de las entidades \parencite[art.~2.2.17.2.1.1, num.~1.1]{colombia2020decreto620}. El marco es un documento técnico; su cumplimiento se vuelve obligatorio por disposición legal para el diseño de los mecanismos de intercambio de información, como se expone en el marco normativo de este capítulo.

El marco de MinTIC señala, además, que la interoperabilidad no es un concepto meramente técnico, y la organiza en cuatro dominios: político-legal, organizacional, semántico y técnico \parencite[p.~33]{mintic2019marco}. La literatura académica describe ese mismo carácter amplio. \textcite{pardo2012interoperability} entienden la interoperabilidad en el gobierno electrónico como un conjunto de capacidades multidimensionales, complementarias y dinámicas que necesitan las redes de organizaciones para compartir información con éxito \pendiente{CONFIRMAR}. Ese carácter amplio también ayuda a explicar sus dificultades: \textcite{scholl2007integration} señalan que la mayoría de los esfuerzos de integración e interoperación enfrentan desafíos y restricciones importantes \pendiente{CONFIRMAR}. En Iberoamérica, \textcite{criado2011interoperabilidad} consultaron a los gobiernos de 16 países \parencite{criado2011interoperabilidad} y encontraron que el reto más mencionado para las iniciativas de interoperabilidad entre Estados era la cultura organizativa, señalada por el 87,5\,\% de los países consultados, y que, en el ámbito tecnológico, el principal problema eran los datos, en particular su dispersión en distintas bases y la ausencia de modelos de datos comunes \pendiente{DATO PENDIENTE: paginación de la revista para Criado et al. (2011)}. Aunque esos resultados corresponden a 2011 y a iniciativas entre países, en este estudio se toman como indicio de que los límites de la interoperabilidad pueden estar tanto en la organización como en los propios datos, y no solo en la tecnología. Por esa razón, los cuatro dominios del marco de MinTIC se usan en el estudio como criterio para clasificar cada documento exigido según lo que tendría que resolverse para dejar de pedirlo: una habilitación legal, un acuerdo entre entidades, un lenguaje común sobre los datos o una conexión técnica.

\subsubsection{Principio de \enquote{solo una vez}}

La idea de que el Estado no pida al ciudadano información que ya posee se conoce en la literatura como el principio de \enquote{solo una vez} (\emph{once-only principle}). \textcite{wimmer2021once} explica, en la introducción de su capítulo, que el principio figura entre los siete principios del plan de acción de administración electrónica 2016-2020 de la Unión Europea, con el fin de hacer el gobierno más eficaz y sencillo y de reducir la carga administrativa, pidiendo a ciudadanos y empresas que entreguen cierta información estándar a las autoridades una sola vez, y que el plan exige a las autoridades tomar medidas para reutilizar internamente esos datos, cuando esté permitido y con respeto de las normas de protección de datos. Su aplicación exige coordinar sistemas de varias entidades: el proyecto europeo con el que se buscó demostrarlo entre países se propuso conectar sistemas de información de varios países e identificar factores impulsores y barreras, como los requisitos de protección e intercambio de datos, los costos de implementación, el aislamiento entre entidades del sector público y los vacíos legales \parencite{krimmer2017exploring} \pendiente{CONFIRMAR}.

En América Latina y el Caribe, \textcite{roseth2018potencial} plantean que la interoperabilidad es lo que permite que el Estado comparta la información ciudadana con la que ya cuenta, de modo que el solicitante no la tenga que volver a presentar, y advierten que la región \enquote{en su mayoría está en una etapa incipiente en cuanto a la interoperabilidad} (p.~120). Según la encuesta BID-GEALC de 2017 que reportan los autores, 13 de 25 países de la región tenían iniciativas de \enquote{pedir una sola vez}, pero solo cuatro de esos 13 sabían cuántas entidades las cumplían \parencite[p.~147]{roseth2018potencial}. De esta manera, que el principio exista en una norma o en un programa no implica que se mida su cumplimiento.

En conjunto, los cuatro conceptos orientan los objetivos propuestos para el estudio, que siguen en revisión. La definición de trámite, la lectura de los documentos como requisitos y su clasificación según quién los expide sustentan la identificación de los documentos que exigen los trámites de la Gobernación del Huila; los conceptos de interoperabilidad, servicio de intercambio de información y Servicios Ciudadanos Digitales, junto con los cuatro dominios del marco de MinTIC, permiten relacionar esos documentos con las bases de datos y los mecanismos de intercambio disponibles y clasificar lo que haría falta para dejar de pedirlos; el principio de \enquote{solo una vez} orienta la propuesta, y la carga administrativa expresa el beneficio que se espera para el ciudadano. La sección siguiente presenta los marcos de interoperabilidad y las normas colombianas que permiten aplicar esos conceptos.

\subsection{Marcos de interoperabilidad y marco normativo}

Esta sección presenta los instrumentos que permiten aplicar los conceptos anteriores al caso: el Marco de interoperabilidad para Gobierno Digital, que ofrece el criterio de clasificación del estudio; su modelo de madurez, del que el estudio toma la forma de las acciones que propone, y las normas colombianas que prohíben exigir documentos disponibles en el Estado y que ordenan la digitalización de trámites y la interoperabilidad, presentadas según su rango y en relación con los instrumentos de planeación de la Gobernación del Huila.

\subsubsection{Los dominios del Marco de interoperabilidad}

El Marco de interoperabilidad para Gobierno Digital parte de que la interoperabilidad no es un concepto meramente técnico, pues involucra retos relacionados con la voluntad política, la formación y apropiación dentro de las entidades, la integración de procesos entre instituciones y la ausencia de un marco legal que faculte a una entidad para intercambiar su información; por esa razón adopta una aproximación por dominios \parencite[p.~33]{mintic2019marco}. El dominio político-legal establece, entre sus lineamientos, contar con instrumentos legales que faciliten el uso o la prestación de los servicios de intercambio de información, fijar la responsabilidad legal de cada parte y garantizar el manejo adecuado de la información confidencial y la protección de datos personales \parencite[p.~34]{mintic2019marco}. El dominio organizacional incluye documentar los procesos, trámites y servicios y concertar la manera de alinearlos para mejorarlos mediante el intercambio de información, adecuarlos para incorporar los servicios de intercambio y establecer las relaciones que se requieren entre las entidades cuando se provee o se consume información para un trámite o servicio \parencite[p.~38]{mintic2019marco}. El dominio semántico busca que, al intercambiar datos, el significado de la información sea exacto y el mismo para todas las partes, y para ello se apoya en el Lenguaje Común de Intercambio de Información, el estándar nacional administrado por MinTIC \parencite[p.~43]{mintic2019marco}. Por último, el dominio técnico se refiere a las aplicaciones e infraestructuras que conectan los sistemas de información a través de los servicios de intercambio, y su primer lineamiento es utilizar los Servicios Ciudadanos Digitales \parencite[p.~47]{mintic2019marco}.

La organización por áreas o dominios no es exclusiva del caso colombiano. Según \textcite{criado2011interoperabilidad}, las Bases para una Estrategia Iberoamericana de Interoperabilidad, aprobadas en 2010, incluyen entre sus apartados uno dedicado a las áreas clave de interoperabilidad técnica, semántica y organizativa \pendiente{DATO PENDIENTE: paginación de la revista; pp. 2-3 del PDF}; en la consulta de esos autores, el Libro Blanco de Interoperabilidad de la Comisión Económica para América Latina y el Caribe y el Marco Europeo de Interoperabilidad figuraban como referentes para el 87,5\,\% y el 81,25\,\% de los países consultados, respectivamente \pendiente{DATO PENDIENTE: paginación de la revista para Criado et al. (2011)}. Para este estudio, los cuatro dominios cumplen la función de criterio de clasificación anunciada en el marco conceptual: a cada documento exigido se le asigna el dominio o los dominios en los que tendría que resolverse algo para que la entidad deje de pedirlo. Esta aplicación es una propuesta del estudio y no un procedimiento que el marco de MinTIC establezca; la revisión de literatura no halló, entre los textos leídos completos, una evaluación crítica de la clasificación por dominios aplicada a trámites, lo que se registra como un límite del criterio.

\subsubsection{El modelo de madurez de interoperabilidad}

El marco incluye un modelo de madurez que permite a cada entidad diagnosticar su avance en la implementación de los lineamientos de cada dominio y definir las acciones para avanzar; el modelo tiene cinco niveles: ausente, inicial, intermedio, consolidado e institucionalizado \parencite[p.~75]{mintic2019marco}. Según el mismo documento, al terminar la evaluación cada acción debe tener una descripción, un responsable y un plazo, y esas acciones deben tenerse en cuenta en los ejercicios de planeación de la entidad \parencite[p.~75]{mintic2019marco}. En las bases consultadas en la revisión de literatura no se halló otra aplicación de este modelo con resultados publicados que la realizada con 46 entidades del orden nacional, que realizaron el diagnóstico entre el último trimestre de 2021 y el primero de 2022: en la calificación general, 10 entidades quedaron en nivel ausente, 7 en inicial, 19 en intermedio, 9 en consolidado y 1 en institucionalizado \parencite[pp.~6, 8]{mintic2022informe}; es decir, 36 de las 46 no habían superado el nivel intermedio (cálculo propio a partir de los conteos, que se citan en lugar de los porcentajes del informe porque dos de estos no coinciden con aquellos). El informe no incluye entidades territoriales, por lo que no ofrece un punto de referencia para la Gobernación del Huila. Como el diagnóstico de madurez es una autoevaluación interna y el estudio trabaja con información pública, el modelo no se aplica aquí para calificar a la Gobernación; de él se adopta la forma de las acciones, de modo que cada acción de la propuesta tenga descripción, responsable y plazo y pueda incorporarse a la planeación de la entidad.

\subsubsection{Marco normativo}

Las normas que sustentan el estudio se presentan según su rango: leyes, decretos con fuerza de ley y decretos reglamentarios. Estos últimos, en el caso del sector de las tecnologías de la información, se incorporan al Decreto Único Reglamentario 1078 de 2015; por ello sus artículos se numeran como 2.2.17.x, en el caso del Decreto 620 de 2020, que subroga el título 17 de ese decreto, y como 2.2.20.x, en el caso del Decreto 088 de 2022, que le adiciona el título 20 \parencites[art.~1]{colombia2020decreto620}[art.~1]{colombia2022decreto088}.

En el nivel de ley, la Ley 962 de 2005 estableció, como principio de la política de racionalización de trámites, que las autoridades no pueden solicitar la presentación de documentos de competencia de otras autoridades \parencite[art.~1, num.~1]{colombia2005ley962}. La Ley 2052 de 2020 fija disposiciones sobre racionalización de trámites para la Rama Ejecutiva del nivel nacional y territorial \parencite[art.~1]{colombia2020ley2052}: ordena que los trámites existentes antes de su vigencia que no puedan realizarse totalmente en línea sigan los plazos y condiciones que determine MinTIC \parencite[art.~6]{colombia2020ley2052}; exige a cada entidad nacional o territorial identificar las cadenas de trámites en las que participa y priorizar la interoperabilidad con las demás entidades que participan en ellas \parencite[art.~8]{colombia2020ley2052}, y obliga a crear, diseñar o adecuar los mecanismos de intercambio de información de los sistemas que soportan los trámites, \enquote{dando cumplimiento al Marcó [sic] de interoperabilidad} y a los lineamientos de vinculación al servicio de interoperabilidad de los Servicios Ciudadanos Digitales \parencite[art.~10]{colombia2020ley2052}. Por esta vía, el marco de MinTIC, que es un documento técnico, se vuelve de cumplimiento obligatorio para los sujetos obligados cuando crean, diseñan o adecuan sus mecanismos de intercambio de información. Por último, la Ley 2195 de 2022 obliga a las entidades del orden nacional, departamental y municipal a implementar Programas de Transparencia y Ética Pública \parencite[art.~31]{colombia2022ley2195}.

En el nivel de los decretos con fuerza de ley, el Decreto Ley 019 de 2012 prohíbe exigir actos administrativos, constancias, certificaciones o documentos \enquote{que ya reposen en la entidad ante la cual se está tramitando la respectiva actuación} y, en su parágrafo, dispone que desde el 1 de enero de 2013 las entidades deben contar con mecanismos para que el solicitante indique la entidad en la que reposan los documentos y esta los requiera directamente, sin perjuicio de que la persona los aporte si así lo prefiere, de modo que no se puede exigir información que repose en los archivos de otra entidad pública \parencite[art.~9]{colombia2012decreto019}. El Decreto Ley 2106 de 2019, expedido con las facultades extraordinarias del artículo 333 de la Ley 1955 de 2019, regula la misma materia en relación con los Servicios Ciudadanos Digitales: obliga a las autoridades a vincular a los mecanismos que disponga la Agencia Nacional Digital los instrumentos y plataformas que contribuyan a esos servicios; dispone que la reutilización de la información integrada al servicio de interoperabilidad se haga bajo las reglas de protección de datos personales y sin necesidad de suscribir acuerdos, convenios o contratos interadministrativos; establece que \enquote{las autoridades no exigirán a los ciudadanos los requisitos o documentos que reposen en bases de datos o sistemas de información que se encuentren integrados en el servicio ciudadano digital de interoperabilidad}; ordena que, cuando las autoridades estén integradas y usando esos servicios, los requisitos verificables por esa vía se actualicen en el Sistema Único de Información de Trámites (SUIT), y dispone que, mientras las autoridades que llevan registros públicos no se integren, habiliten su consulta gratuita y en línea a las demás autoridades, de modo que el ciudadano quede eximido de aportar el certificado \parencite[art.~10]{colombia2019decreto2106}.

En el nivel reglamentario, el Decreto 620 de 2020 obliga a las entidades a incluir \enquote{los mecanismos de interoperabilidad necesarios que permitan hacer más ágiles y eficientes los trámites y servicios evitando solicitar información a ciudadanos y empresas a la que puedan acceder, consultar o solicitar a otra entidad} \parencite[art.~2.2.17.4.7, num.~8]{colombia2020decreto620}. El Decreto 088 de 2022, que reglamenta los artículos 3, 5 y 6 de la Ley 2052 de 2020, adopta como definición normativa de interoperabilidad la misma del marco de MinTIC, con la diferencia de que se refiere a la entrega de \enquote{servicios digitales} \parencite[art.~2.2.20.3, num.~8]{colombia2022decreto088}, y obliga a las autoridades a adoptar y aplicar la guía de su Anexo 1 \parencite[art.~2.2.20.4]{colombia2022decreto088}. Esa guía contiene la disposición más cercana a lo que hace este estudio: con base en los requisitos legales del trámite, se deben identificar los documentos e información que se piden al usuario, verificar cuáles ya tiene disponibles la entidad e identificar cuáles corresponden a otras entidades según su competencia, para luego coordinar con estas la habilitación de los servicios de intercambio, siguiendo el Lenguaje Común de Intercambio de Información y el Marco de interoperabilidad, \enquote{con el propósito de evitar pedir nuevamente la información con la que ya se cuenta del usuario} \parencite[Anexo 1, num.~2.3.3.1.5]{colombia2022decreto088}. La misma guía exige a las entidades un plan de estandarización de sus servicios de intercambio de información ajustado al Lenguaje Común \parencite[Anexo 1, num.~2.4.2]{colombia2022decreto088}, y establece que las iniciativas de digitalización o automatización de trámites deben incorporarse en el Plan Estratégico de Tecnologías de la Información (PETI) y, cuando generan un beneficio directo al ciudadano, registrarse además como acciones de racionalización en el plan anticorrupción \parencite[Anexo 1, num.~2.3.2]{colombia2022decreto088}. Para las gobernaciones, el decreto fija que el 30\,\% de los trámites debe estar digitalizado hasta mayo de 2028, el 60\,\% hasta julio de 2031 y el total hasta marzo de 2034, entendida la digitalización como el uso de medios digitales en la gestión interna de los trámites \parencite[arts. 2.2.20.3, num.~3, y 2.2.20.7]{colombia2022decreto088}; para las entidades territoriales, esos plazos están sujetos a las condiciones de conectividad, infraestructura y tecnologías requeridas y a la disponibilidad de presupuesto, y pueden solicitar de manera motivada su ampliación \parencite[art.~2.2.20.6, num.~3]{colombia2022decreto088}. Finalmente, el Decreto 1122 de 2024, mediante su artículo 2, adiciona al Decreto 1081 de 2015 la regla según la cual toda mención al Plan Anticorrupción y de Atención al Ciudadano se entiende referida al Programa de Transparencia y Ética Pública \parencite[art.~2]{colombia2024decreto1122}.

En cuanto a la vigencia, la Ley 962 de 2005, la Ley 2052 de 2020, la Ley 2195 de 2022, el Decreto Ley 019 de 2012, el Decreto Ley 2106 de 2019, el Decreto 088 de 2022 y el Decreto 1122 de 2024 figuran como vigentes en el Sistema Único de Información Normativa (SUIN-Juriscol) del Ministerio de Justicia, que informa el estado de cada norma en su conjunto; para el Decreto 620 de 2020, la ficha consultada no registra estado de vigencia \pendiente{DATO PENDIENTE: vigencia del Decreto 620 de 2020}. Respecto del artículo 10 del Decreto Ley 2106 de 2019, la Corte Constitucional, en la Sentencia C-023 de 2021, se inhibió de pronunciarse de fondo sobre los artículos 9 a 13 de ese decreto, y la compilación anotada de Avance Jurídico, que no es una fuente oficial, no registra modificaciones a ese artículo \pendiente{CONFIRMAR: cotejar el texto completo de la sentencia}.

De este conjunto de normas se deriva qué regla aplica a cada categoría de documento definida en el marco conceptual. Para los documentos que reposan en la propia Gobernación, la prohibición de exigirlos es directa y no depende de ninguna integración tecnológica \parencite[art.~9]{colombia2012decreto019}. Para los documentos de competencia de otras autoridades, la prohibición también existe \parencites[art.~1, num.~1]{colombia2005ley962}[art.~9]{colombia2012decreto019}, pero el mecanismo varía: si la entidad que expide el documento está integrada al servicio de interoperabilidad, el documento no puede exigirse y la información está disponible por esa vía sin necesidad de convenios; si se trata de un registro público aún no integrado, debe consultarse en línea, y en los demás casos el solicitante puede indicar la entidad en la que reposa el documento para que la Gobernación lo requiera \parencites[art.~9]{colombia2012decreto019}[art.~10]{colombia2019decreto2106}. Para la información que solo el ciudadano puede aportar, las normas revisadas no contienen una regla específica; en este estudio se trata como la categoría que permanece como requisito, sin perjuicio de que pueda simplificarse por otras vías de racionalización. Esta lectura conjunta es una interpretación del estudio, y su aplicación a cada documento se precisa en la metodología.

La Gobernación del Huila ya reconoce este marco en sus propios instrumentos. El marco normativo del PETI 2026 incluye, entre otras, la Ley 2052 de 2020, el Decreto Ley 019 de 2012, el Decreto 2106 de 2019, así denominado en ese documento, y el Decreto 088 de 2022 \parencite[pp.~6-8]{gobernaciondelhuila2026peti}; el Programa de Transparencia y Ética Pública 2026-2027 incluye la racionalización de trámites como iniciativa adicional \parencite[p.~133]{gobernaciondelhuila2026ptep} y su diagnóstico identifica \enquote{una necesidad de profundizar en la interoperabilidad de bases de datos, ya que persiste la solicitud de documentos que la entidad ya posee o que pueden consultarse en línea} \parencite[p.~134]{gobernaciondelhuila2026ptep}, y el Plan de Desarrollo Departamental 2024-2027 incluye, dentro del sector Gobierno Territorial, dos estrategias relacionadas: en el apartado de gobernanza institucional, revisar y ajustar procesos y procedimientos para disminuir trámites y mejorar la eficiencia administrativa en el acceso del ciudadano \parencite[p.~187]{gobernaciondelhuila2024ordenanza}, y, en el componente \enquote{Huila financia su futuro}, mejorar los servicios tecnológicos mediante sistemas de información internos con interoperabilidad entre procesos y con otros niveles de gobierno o entidades, entre ellas la banca, el registro y la cámara de comercio \parencite[p.~190]{gobernaciondelhuila2024ordenanza}. Estas menciones muestran que la interoperabilidad figura en la planeación departamental; la relación del estudio con un producto o meta específicos del Plan se precisa en la justificación \pendiente{DATO PENDIENTE: producto o meta del Programa 4599 al que se ancla el estudio}. Queda por establecer si la Gobernación está integrada al servicio ciudadano digital de interoperabilidad, condición de la que dependen el mecanismo de consulta de sus propias bases y la obligación del parágrafo 2 del artículo 10 del Decreto Ley 2106 de 2019 \pendiente{DATO PENDIENTE: integración de la Gobernación del Huila a los Servicios Ciudadanos Digitales}. La sección siguiente revisa los antecedentes y el estado del arte.

\subsection{Antecedentes y estado del arte}

Esta sección revisa estudios y experiencias sobre la aplicación del principio de \enquote{solo una vez} y de la interoperabilidad para reducir la información que se pide al ciudadano. Se presentan los antecedentes que se leyeron en texto completo y, para la evidencia que matiza los beneficios esperados, también trabajos leídos en resumen, de los que solo se recoge el hallazgo general. Se organiza de lo internacional a lo local: Europa, Iberoamérica, Colombia en el nivel nacional y Colombia en el nivel territorial; al final se reúne la evidencia que matiza y se señala el vacío.

\subsubsection{Europa}

La experiencia europea es la más documentada. \textcite{mamrot2021implementation} compararon la forma en que los países de la Unión Europea y del Espacio Económico Europeo aplican el principio en el intercambio de datos de empresas, y concluyeron que en la mayoría está incorporado tanto en la legislación nacional como en las estrategias de gobierno electrónico; sin embargo, advierten que tener la norma \enquote{no equivale} al funcionamiento práctico del principio, y que algunos países solo cuentan con soluciones limitadas a ciertos servicios o registros (p.~17). Entre los casos que describen, la ley de Austria dispone que, cuando sea técnicamente posible, no se le pida al ciudadano probar datos que ya existen en un registro electrónico del sector público, y que las entidades los consulten directamente en las bases de datos (p.~18); en Estonia, las instituciones públicas intercambian información a través del sistema X-Road, cuyo uso impone la legislación (pp.~11-12).

\textcite{wimmer2021once} informa que en el proyecto SCOOP4C se analizaron 57 casos de aplicación del principio y 34 elementos que la facilitan, y que se identificaron once condiciones clave, entre ellas el compromiso político, el marco legal, el compromiso organizacional con procesos compartidos entre entidades, la interoperabilidad semántica mediante estándares comunes de intercambio, la infraestructura técnica segura y la protección de datos (p.~77). La autora señala que las barreras y las condiciones habilitantes son conceptos ligados, porque una barrera indica la ausencia de una condición habilitante (p.~77), y que, aunque las políticas europeas promueven ampliamente la interoperabilidad, la transformación digital en la práctica va considerablemente rezagada respecto de esas visiones (p.~78). Por su parte, \textcite{leosk2021drivers}, que estudiaron la aplicación transfronteriza del principio en la Unión Europea, encontraron que la coordinación entre distintos niveles de gobierno y entre países sigue siendo una barrera importante (p.~38) y que, en varios de los casos estudiados, las conexiones entre quienes producen y quienes consumen datos se establecieron punto a punto, de forma bilateral, lo que optimizó el intercambio entre unos pocos sistemas pero produjo falta de interoperabilidad (p.~54).

\subsubsection{Iberoamérica}

Los antecedentes de la región, ya presentados en el marco conceptual \parencites{criado2011interoperabilidad}{roseth2018potencial}, son diagnósticos de nivel nacional. A ellos se suma que, en la consulta de \textcite{criado2011interoperabilidad} a 16 gobiernos iberoamericanos, el 75\,\% de los países declaraba contar con una estrategia de interoperabilidad y los trámites generales figuraban entre las áreas con proyectos de interoperabilidad en el 62,5\,\% de los países \pendiente{paginación de la revista para Criado et al. (2011)}. Ninguno de estos trabajos examina entidades subnacionales.

\subsubsection{Colombia: nivel nacional}

\textcite{moralesmoreno2015interoperability} describen, con base en un diagnóstico que retoman de fuentes oficiales, el punto de partida de la estrategia colombiana de interoperabilidad como servicio: las entidades públicas operaban como silos y, cuando un trámite requería información de otra entidad, era el ciudadano quien terminaba conectando los servicios al desplazarse físicamente de una entidad a otra (p.~66). Entre las barreras que la iniciativa debía atender señalan el interés de las entidades en publicar sus datos, las capacidades de desarrollo de software, la capacidad de infraestructura y la calidad de los datos (p.~71). La medición del modelo de madurez descrita en el marco normativo tampoco incluye entidades territoriales. En el nivel de una entidad, \textcite{molina2025propuesta} propuso una estrategia de interoperabilidad para los servicios ciudadanos de una organización del sector defensa, basada en una plataforma en la nube, e identificó como desafíos la infraestructura tecnológica y la necesidad de un cambio cultural en las organizaciones (p.~4); anticipa, además, que la capacitación del personal y la adecuación de los sistemas existentes podrían obstaculizar su implementación (p.~67).

\subsubsection{Colombia: nivel territorial}

Los estudios sobre entidades territoriales son escasos. \textcite{marquezpinzon2023modelo} diagnosticaron los trámites de cuatro alcaldías de la provincia de Ocaña y propusieron un modelo de gestión tecnológica basado en COBIT 2019. En sus conclusiones reportan, con base en las mediciones del FURAG y sin precisar el año, una diferencia de 25,2 puntos en el índice de gobierno digital y de 26,4 en el índice de desempeño institucional entre las entidades del orden nacional y las del orden territorial, y señalan que solo el municipio de El Carmen había incluido por completo sus trámites en el Sistema Único de Información de Trámites (p.~72). En otro sector, \textcite{pinzonarana2025caracterizacion} describen cómo el Sistema de Administración del Territorio se reorganizó a partir de la interoperabilidad entre catastro y registro (p.~6), afirman que en los trámites catastrales el uso de sistemas interoperables ha reducido requisitos y tiempos (p.~19) y explican que la estandarización de la información fue el primer paso para lograrlo, porque cada entidad producía su propia versión de un mismo objeto territorial (p.~20).

\subsubsection{Evidencia que matiza}

Otros trabajos matizan la expectativa de que la interoperabilidad reduzca por sí sola la carga del ciudadano. En el proyecto TOOP, \textcite{timan2021measuring} encontraron que se sabe poco del impacto del principio sobre las empresas, en particular del impacto transfronterizo, y que sus beneficios esperados probablemente aparecen a largo plazo, lo que dificulta identificar impactos de corto plazo (p.~208). En el nivel local, el intercambio electrónico de información entre autoridades depende de factores del entorno, organizacionales, de procesos y tecnológicos, y no solo técnicos \parencite{bigdeli2013electronic}; en la misma línea, a partir de la experiencia de Estonia, \textcite{puura2026advancing} sostienen que los principales obstáculos para el intercambio de datos son organizacionales, políticos y legales más que tecnológicos. Compartir datos personales entre entidades tiene, además, un costo en protección de datos y privacidad: los ciudadanos lo perciben así, sobre todo con datos sensibles \parencite{akkaya2018towards}, y el principio de \enquote{solo una vez} entra en tensión con el principio de limitación de la finalidad de la protección de datos \parencite{mikiver2023gdpr}.

\subsubsection{Síntesis y vacío}

De los antecedentes se desprenden tres puntos. Primero, que la existencia de una norma o de una política no garantiza que el principio de \enquote{solo una vez} se aplique en la práctica \parencites[p.~17]{mamrot2021implementation}[p.~78]{wimmer2021once}. Segundo, que los obstáculos no son solo técnicos: en los estudios sobre interoperabilidad entre países, la coordinación entre niveles de gobierno y la cultura organizativa aparecen como obstáculos importantes \parencites{criado2011interoperabilidad}[p.~38]{leosk2021drivers}, y en el nivel local pesan factores organizacionales y de procesos \parencite{bigdeli2013electronic}; en Colombia, las barreras previstas incluyen además la disposición de las entidades a publicar sus datos y la calidad de esos datos \parencite[p.~71]{moralesmoreno2015interoperability}. Tercero, que los beneficios del principio están poco medidos y la protección de datos personales limita su alcance \parencites[p.~208]{timan2021measuring}{mikiver2023gdpr}. En el caso colombiano, \textcite[p.~72]{marquezpinzon2023modelo} reportan, con mediciones del FURAG cuyo año no precisan, que las entidades del orden territorial obtuvieron puntajes más bajos que las del orden nacional en el índice de gobierno digital. En las bases consultadas (Consensus, SciSpace y búsqueda dirigida a SciELO, Redalyc, Dialnet y Google Académico, en octubre de 2026) no se halló ningún estudio sobre la interoperabilidad o la reducción de requisitos en los trámites de una gobernación colombiana. Tampoco se halló, en Colombia, un estudio que parta de los documentos exigidos en el conjunto de trámites de una entidad para identificar cuáles podrían obtenerse por interoperabilidad: los antecedentes más cercanos tratan un sector, como el catastro y el registro \parencite{pinzonarana2025caracterizacion}, o entidades del orden nacional \parencites{moralesmoreno2015interoperability}{molina2025propuesta} \pendiente{incorporar y cotejar con página el Marco de interoperabilidad de la Alcaldía Distrital de Cartagena de Indias, que según la revisión menciona los trámites afectados por los intercambios pero no la reducción de los documentos pedidos al ciudadano}. La búsqueda no completó el rastreo de las referencias citadas por los antecedentes, limitación que se declara en la metodología. Ese es el vacío que aborda este estudio.

Para el estudio, estos antecedentes tienen tres usos, que son una propuesta de esta investigación y no un hallazgo de las fuentes. La estandarización de la información como paso previo a la interoperabilidad \parencite[p.~20]{pinzonarana2025caracterizacion} orienta la revisión del dominio semántico al relacionar cada documento con las bases de datos que lo respaldan. Los efectos de las conexiones bilaterales entre sistemas \parencite[p.~54]{leosk2021drivers} y las barreras previstas en Colombia \parencite[p.~71]{moralesmoreno2015interoperability} sirven como criterios al diseñar la propuesta. Por último, los obstáculos que \textcite[p.~67]{molina2025propuesta} anticipa para una propuesta similar, junto con la evidencia que matiza los beneficios, se usan como criterios para valorar la factibilidad de la propuesta \pendiente{ajustar las referencias a los objetivos específicos cuando se aprueben}.
